\chapter{Introduction}

\section{Problem Description and Motivation}
\label{sec:problem}

The exponential growth of emerging technologies---including artificial intelligence, machine learning, distributed systems, and advanced computing---has produced an unprecedented volume of technical knowledge distributed across heterogeneous platforms. Scholarly repositories such as ArXiv provide structured, peer-reviewed or preprint research outputs contributing to theoretical and methodological advancement. Community-driven platforms and industry discussion forums such as Hacker News reflect near-real-time practitioner experiences and emerging tool ecosystems. However, the fragmentation of these knowledge sources creates an information landscape that is increasingly difficult for technology professionals to synthesize effectively and in a timely manner.

Purely generative large language model (LLM) assistants are insufficient for emerging technology intelligence tasks due to their static knowledge cutoff, susceptibility to hallucinated responses, lack of verifiable evidence grounding, and inability to integrate with continuously updating data streams (Ji et al., 2023). Although Retrieval-Augmented Generation (RAG) architectures partially address these limitations by incorporating external context, the relative benefit of augmenting vector-only retrieval with lexical and recency signals remains under-evaluated in cost-constrained, topically focused settings. Moreover, baseline vector-only systems typically lack temporal relevance signals, structured metadata constraints, and recency-aware ranking mechanisms, which are important for accurately capturing rapidly evolving technology trends.

A further gap exists at the systems level: many RAG research prototypes focus primarily on retrieval accuracy without addressing production-readiness---deterministic pipeline design, idempotent ingestion, monitoring and observability, deployment automation, and cloud cost governance are frequently neglected, limiting their applicability in real-world operational environments.

The central engineering and research challenge addressed in this project is therefore:

\begin{quote}
\textbf{\textit{How can a scalable, continuously updating hybrid retrieval system be designed and deployed to synthesize emerging technology trends from heterogeneous scholarly and industry sources while satisfying latency and cost constraints?}}
\end{quote}

This project is motivated by three converging factors. First, technology professionals increasingly require timely and citation-grounded intelligence that synthesizes insights from both scholarly research and rapidly evolving industry discussions. Second, mature serverless cloud infrastructures now enable cost-effective deployment of production-grade data and machine learning pipelines. Third, a gap remains in the Retrieval-Augmented Generation (RAG) literature between retrieval-focused research prototypes and operationally robust systems capable of continuous ingestion, evaluation, and deployment in real-world environments.

\section{Research Questions}
\label{sec:rq}

To address the central problem, this study investigates the following research questions:

\begin{enumerate}[leftmargin=0.5in]
  \item \textbf{Data Engineering.} How can a cloud-native pipeline reliably ingest heterogeneous scholarly and industry data in a deterministic, idempotent manner?

  \item \textbf{Retrieval Optimization.} Does hybrid retrieval (metadata filtering, BM25, and Reciprocal Rank Fusion) improve precision and faithfulness over baseline vector-only retrieval?

  \item \textbf{Operational Trade-offs.} What trade-offs arise between retrieval precision, p95 latency, and cost per query?

  \item \textbf{MLOps Integration.} How can deployment automation, monitoring, and cost governance ensure the reliability of a production-grade RAG system?
\end{enumerate}

\section{Objectives of the Study}
\label{sec:objectives}

The primary objective is to design, implement, and evaluate a hybrid Retrieval-Augmented Generation system capable of synthesizing real-time emerging technology intelligence within a scalable, cost-aware cloud-native architecture. The specific objectives are:

\begin{enumerate}[leftmargin=0.5in]
  \item To construct a cloud-native data engineering pipeline for batch and near-real-time ingestion from five heterogeneous data sources.
  \item To implement and compare baseline vector-only retrieval against hybrid retrieval with metadata filtering, full-corpus BM25, and Reciprocal Rank Fusion (RRF).
  \item To evaluate system performance across context precision, faithfulness, p95 retrieval latency, and cost per query.
  \item To deploy and monitor the system using serverless AWS infrastructure consistent with MLOps best practices (Sculley et al., 2015; Breck et al., 2017).
\end{enumerate}

\section{Significance of the Study}
\label{sec:significance}

This project contributes at the intersection of data engineering, natural language processing, and MLOps in three ways. First, it provides controlled empirical evidence on whether hybrid retrieval (BM25 + RRF) improves over vector-only baselines within a topically homogeneous, cost-constrained setting---a configuration common in practice but under-represented in the literature. Second, it demonstrates that parameter-free Reciprocal Rank Fusion achieves comparable retrieval quality to tuned weighted reranking while eliminating dataset-dependent weight sensitivity and reducing tail latency by 87\%. Third, the end-to-end system integrates hybrid retrieval, drift detection, hallucination verification, and cost governance within a single deployable architecture on the AWS Free Tier, providing a practical blueprint for student and small-team RAG deployments without reliance on expensive cross-encoder rerankers or dedicated vector database infrastructure.

\section{Alignment Between Research Questions and Evaluation}
\label{sec:rq_eval_map}

Table~\ref{tab:rq_map} provides an explicit mapping from each research question to the corresponding evaluation metric and experimental configuration used to answer it.

\begin{table}[H]
\caption{Mapping of Research Questions to Evaluation Metrics and Methodology.}
\label{tab:rq_map}
\centering
\begin{tabular}{p{0.5cm}p{4.0cm}p{3.5cm}p{4.0cm}}
\toprule
\textbf{RQ} & \textbf{Question (abbreviated)} & \textbf{Primary Metric} & \textbf{Evaluation Approach} \\
\midrule
1 & Cloud-native deterministic ingestion & Pipeline idempotency; SHA-256 deduplication pass rate & Unit tests (24 ingestion tests); repeated pipeline runs \\[4pt]
2 & Hybrid vs.\ baseline retrieval precision & Context precision (RAGAS); faithfulness $\geq 0.80$; citation grounding $> 0.80$ & Paired comparison over 50 queries; Wilcoxon signed-rank test; Cohen's $d$; bootstrap 95\% CI \\[4pt]
3 & Latency and cost trade-offs & p95 latency $< 2$s (retrieval stage); cost $<$ USD~15/month & Per-query timing; token usage profiling; cost projection \\[4pt]
4 & MLOps integration & Drift detection coverage; CI/CD pass rate; health-check uptime & Drift simulation runs; CloudWatch alarm verification \\
\bottomrule
\end{tabular}
\end{table}

\medskip
\noindent The study is structured as follows: Chapter~2 reviews the relevant literature; Chapter~3 presents the methodology and hypotheses; Chapter~4 describes the system architecture; Chapter~5 defines the evaluation framework; Chapter~6 documents implementation progress; and Chapter~7 outlines remaining work and limitations. The following chapter reviews the technical foundations underpinning these design decisions.
